\chapter{时序特征提取与匹配}
\label{chap:methods}

\section{样本构建}

\subsection{样本划分}

以原始视频为单位划分调参与评估样本。同一原片的重编码、裁切或分辨率变体归入同一组。检索集同时保留重复片段和视觉内容相近的非重复片段。

\begin{table}[htbp]
  \centering
  \caption{片段类别与评估用途}
  \label{tab:groups}
  \begin{tabular}{ll}
    \toprule
    片段类别 & 评估用途 \\
    \midrule
    同源重编码片段 & 检查对编码变化的稳定性 \\
    同源裁切片段 & 检查对画面裁切的稳定性 \\
    相似场景的异源片段 & 检查误匹配 \\
    \bottomrule
  \end{tabular}
\end{table}

\section{特征匹配}

对长度为 \(T\) 的窗口，将各帧特征 \(\mathbf{f}_t\) 求平均，再将窗口描述符归一化。
\begin{equation}
  \mathbf{v} = \frac{1}{T}\sum_{t=1}^{T}\mathbf{f}_t .
  \label{eq:pooling}
\end{equation}
按归一化描述符的内积对候选片段排序。特征范数为零的窗口单独标记，并从匹配评分中排除。图~\ref{fig:workflow}~与表~\ref{tab:groups}~分别给出处理流程和样本用途。

\begin{figure}[htbp]
  \centering
  \fbox{帧特征提取}\enspace\(\longrightarrow\)\enspace
  \fbox{窗口池化}\enspace\(\longrightarrow\)\enspace
  \fbox{相似度排序}
  \caption{短时片段匹配流程}
  \label{fig:workflow}
\end{figure}
